\section{落地指南：从开源 Agent 到团队生产系统}

如果把 OpenClaw 视作“起点代码库”，团队还需要走完一条产品工程链。下面给出可执行路线图，帮助把实验项目升级为可运营系统。

\subsection{阶段一：单人可用（个人生产力）}

目标是让一个用户在可控风险内持续得到价值。重点不是并发，而是闭环完整。

\begin{knowledgebox}{阶段一必做}
\begin{itemize}
    \item 定义 2-3 个稳定任务场景；
    \item 约束工具权限，禁止高风险自动执行；
    \item 建立最小日志与回放能力；
    \item 每周复盘失败样例并修正规则。
\end{itemize}
\end{knowledgebox}

\subsection{阶段二：小团队共用（协作与治理）}

团队接入后，问题从“能不能做”变成“谁对结果负责”。需要引入角色、审计与配置管理。

\begin{importantbox}{阶段二的核心改造}
\begin{enumerate}
    \item 多用户权限模型（Owner/Operator/Viewer）；
    \item 环境隔离（dev/staging/prod）；
    \item 配置版本化与审批流；
    \item 高风险动作双人确认机制。
\end{enumerate}
\end{importantbox}

\subsection{阶段三：规模化运营（稳定性与成本）}

进入规模化后，系统需同时优化稳定性、成本和延迟。此时最容易犯的错误是“为了省成本过度降级模型”，导致整体体验崩塌。

\begin{warningbox}{规模化阶段的三类反模式}
\begin{itemize}
    \item 只看 token 成本，不看返工与事故成本；
    \item 把 prompt 调优当成主要优化手段，忽视系统瓶颈；
    \item 缺乏回归测试，导致修一处坏三处。
\end{itemize}
\end{warningbox}

\subsection{评测框架：不要只看 benchmark}

可将评测分为四层：
\begin{itemize}
    \item \textbf{Task Success}：任务完成率、人工接管率；
    \item \textbf{Reliability}：重试率、崩溃率、恢复时间；
    \item \textbf{Safety}：越权动作数、误触发高风险操作数；
    \item \textbf{Business}：节省工时、用户留存、净收益。
\end{itemize}

这四层指标应同时监控，否则“局部优化”会变成系统性退化。

\subsection{人才结构：Agent 时代团队如何配置}

访谈里“builder identity”的观点对应到团队配置，通常需要：
\begin{itemize}
    \item 产品型工程师：负责任务建模和价值闭环；
    \item 平台工程师：负责 runtime、权限、审计体系；
    \item 安全工程师：负责对抗样例、攻防演练与事件响应；
    \item 领域专家：负责规则约束与结果验收。
\end{itemize}

\begin{knowledgebox}{组织层面的关键转变}
团队不再按“前后端边界”严格分工，而更像围绕“任务闭环”协作。谁最懂任务，谁就应参与系统约束设计。
\end{knowledgebox}

\subsection{本章小结}
从开源 demo 到生产系统，真正的工作量在“治理层”而非“首版功能层”。OpenClaw 给出的最大价值，是提供了一条可观察、可学习、可再工程化的路线。

